Furthermore, our approach differs from VectorFusion~\cite{jain2022vectorfusion} mainly in two aspects: 1) Initialization: VectorFusion randomly initializes the rasterizer (i.e., the location of the control points), while our approach utilizes the attention layer of LDM, guided by text prompts, for initialization. This significantly improve both efficiency and syntehsis quality. 2) Optimization: While VectorFusion solely optimizes the position of control points, our method also optimize the opacity of stroke. This contributes to enhance the visual quality of the synthesized sketches.
DiffSketcher produces much cleaner vector sketches than baselines because we incorporate a generative prior of image appearance.

Our method, DiffSketcher, focuses on generating vector sketches from text input. It differs from VectorFusion, which aims to generate a wider range of vector graphics. DiffSketcher can easily extend its functionality to generate other vector graphics without altering primitives. The approach of DiffSketcher follows a distinct pipeline compared to VectorFusion. We show the results of our approach in generating general vector graphics in Fig.~\ref{sec:compare_vs_VF} . The differences between the two methods are mainly in initialization and optimization. Our approach utilizes the attention layer of LDM guided by text prompts for initialization, while VectorFusion randomly initializes the rasterizer. In optimization, DiffSketcher optimizes both the position of control points and the opacity of stroke, while VectorFusion solely optimizes the position of control points. We present a performance comparison between these two methods in terms of vector sketch generation in Fig.~\ref{fig:VF_vs_ours_scratch}.
